\documentclass[12pt,a4paper]{article}
\usepackage[utf8]{inputenc}
\usepackage[portuguese]{babel}
\usepackage{amsmath}
\usepackage{graphicx}
\usepackage{hyperref}
\usepackage{cite}
\usepackage{booktabs}

\title{Estudo Analítico de Diodos Semicondutores \\ \large Baseado em Sedra/Smith e Aplicações Modernas}
\author{Cleber da Silva Alves}
\date{\today}

\begin{document}

\maketitle

\section{Introdução aos Semicondutores}
O diodo é o dispositivo semicondutor mais simples, consistindo em uma junção PN. Segundo Sedra e Smith \cite{sedra}, o comportamento do diodo é fundamental para entender dispositivos mais complexos como transistores BJT e MOSFET.

\section{Fundamentação Teórica e Equações}
\subsection{A Equação de Shockley}
A característica corrente-tensão ($I-V$) de um diodo ideal é descrita pela equação de Shockley:
\begin{equation}
I_D = I_S \left( e^{\frac{V_D}{n V_T}} - 1 \right)
\end{equation}
Onde:
\begin{itemize}
    \item $I_S$: Corrente de saturação reversa (típica de $10^{-15}$ A a $10^{-12}$ A).
    \item $V_D$: Tensão aplicada no diodo.
    \item $V_T$: Tensão térmica, $V_T = \frac{kT}{q} \approx 25 \text{ mV}$ a 300 K.
    \item $n$: Fator de emissão ou fator de idealidade ($1 \leq n \leq 2$).
\end{itemize}

\subsection{Modelos de Análise}
Para análise de circuitos, utilizamos três modelos principais:
\begin{enumerate}
    \item \textbf{Modelo Exponencial:} Mais preciso, exige métodos iterativos ou gráficos.
    \item \textbf{Modelo de Queda de Tensão Constante:} Assume $V_D = 0,7\text{V}$ quando conduzindo.
    \item \textbf{Modelo de Diodo Ideal:} Curto-circuito em polarização direta e circuito aberto em reversa.
\end{enumerate}

Segundo Malvino \cite{malvino}, a compreensão das aproximações do diodo (primeira, segunda e terceira aproximações) é vital para o diagnóstico de defeitos em circuitos práticos, permitindo uma análise rápida antes do refinamento matemático.

\section{Circuitos com Diodos}
\subsection{Retificadores e Detectores de Pico}
O retificador de meia onda converte AC em DC pulsante. Com um capacitor em paralelo à carga, obtemos um detector de pico. A tensão de ondulação (ripple) $\Delta V_r$ é dada por:
\begin{equation}
\Delta V_r = \frac{V_p}{f R C}
\end{equation}
Onde $V_p$ é a tensão de pico e $f$ a frequência da rede.

\subsection{Diodo Zener e Regulação}
O diodo Zener opera na região de ruptura. O modelo linear inclui a tensão nominal $V_Z$ e a resistência dinâmica $r_z$:
\begin{equation}
V_Z = V_{Z0} + r_z I_Z
\end{equation}

\section{Avanços e Tendências (2025)}
Pesquisas recentes destacam o uso de novos materiais para superar as limitações do silício:
\begin{itemize}
    \item \textbf{Diodos de Nitreto de Gálio (GaN):} Oferecem maior eficiência em fontes chaveadas e carregadores rápidos \cite{gan2025}.
    \item \textbf{Diodos de Carbeto de Silício (SiC):} Cruciais para eletrônica de potência em veículos elétricos devido à alta estabilidade térmica \cite{sic2025}.
    \item \textbf{Perovskitas em LEDs:} Avanços em 2025 permitem LEDs azuis de alta eficiência com baixa queda de performance \cite{perovskite2025}.
\end{itemize}

\begin{thebibliography}{9}
\bibitem{sedra} SEDRA, A. S.; SMITH, K. C. \textbf{Microelectronic Circuits}. 7th ed. Oxford University Press, 2014.
\bibitem{gan2025} PIETRZAK, A. Development of high-efficiency and thermally stable pump diodes. \textit{SPIE Digital Library}, 2026.
\bibitem{sic2025} TechInsights. Five Key Trends for Power Semiconductors in 2025.
\bibitem{perovskite2025} TANG, J. et al. High-Efficiency Blue Perovskite Light-Emitting Diodes. \textit{Advanced Materials}, 2026.
\bibitem{malvino} MALVINO, A. P.; BATES, D. J. \textbf{Eletrônica}. 8ª ed. McGraw-Hill, 2016.
\end{thebibliography}

\end{document}
